\subsection*{Transpilation and Compilation}

The Croqtile compilation process typically involves three major steps: \textbf{Pre-processing}, \textbf{Transpilation}, and \textbf{Target Compilation}. To better understand how different parts of a Croqtile program work together, the full compilation workflow is illustrated below:

\begin{figure}[htbp]\centering
\includegraphics[width=0.86\linewidth,height=0.42\textheight,keepaspectratio]{f14-compile-workflow.drawio.png}
\end{figure}

As shown in the figure, immediately after \emph{pre-processing}, Croqtile \emph{transpiles} (\emph{source-to-source} compilation) the \emph{tileflow program} into target code, leaving the user-provided \emph{host program} and \emph{device program} unchanged. The \emph{tileflow program} is \emph{transpiled} into host and device source code, which you refer to as \emph{croqtile-generated code}. The compiler then combines the user-provided code and \emph{croqtile-generated} code to perform the \emph{target compilation} process. This process can generate various outputs, such as transpiled source code, workscripts, target modules, target assembly, and executables.

The Croqtile compiler functions as an end-to-end compiler, with the key step being the transpilation of the \emph{tileflow program} into \emph{croqtile-generated code}.

One notable feature of Croqtile compilation is its support for both:

\begin{itemize}[leftmargin=1.6em]
\item \textbf{Single Source Compilation Model}: Similar to , where the target compiler allows device and host programs to be in a single source file for target compilation.
\item \textbf{Separate Source Compilation Model}: Similar to , where host and device code must be compiled separately.
\end{itemize}

The code shown above naturally supports the \emph{Single Source Compilation Model}. However, to support the \emph{Separate Source Compilation Model}, Croqtile requires wrapping the Device Program with the \texttt{\_\_cok\_\_} block, as shown below:

\begin{lstlisting}[language=croqtile,caption={Transpilation and Compilation}]
__cok__ {
  void device_function(...) { ... }
} // end of __cok__
__co__ void croqtile_function(...) { ... }
void foo() { ... }
\end{lstlisting}

This is the code structure for a  program. The \texttt{\_\allowbreak{}\_\allowbreak{}cok\_\allowbreak{}\_\allowbreak{} \{\}} wrapper enables the Croqtile compiler to handle user-provided device code properly. It helps Croqtile separate device and host code from a single Croqtile source file for different compilation processes. Therefore, do not be surprised if you encounter \texttt{\_\_cok\_\_} in certain Croqtile code; it is necessary for integrating support for the \emph{Separate Source Compilation Model}.
